\section{Discussion}
While FastWrite shows promise, there are several open challenges and limitations worth noting.
\subsection{Limitations}

\textbf{Context Awareness:} The current AI model processes text in small segments, such as paragraphs or sections, which can overlook cross-sectional inconsistencies.
For instance, a variable defined in Section 3 might be used differently in Section 2, and the AI may not detect this discrepancy.
Expanding the context window to address this issue is a priority, but it increases latency and computational costs.

\textbf{PDF Compilation:} FastWrite focuses on the writing experience.
For final PDF compilation, users rely on external tools like `latexmk` or Overleaf.
We integrate a local PDF previewer using the native TeX engines installed by user.

\subsection{Future Work}
\begin{itemize}
    \item \textbf{Local LLM Support:} To enhance privacy and offline capability, we are experimenting with quantized models, such as Llama-3-8B, running locally via Ollama.
    \item \textbf{Collaboration:} We will implement CRDT-based real-time collaboration, enabling multiple authors to edit the same local file through P2P connections.
    \item \textbf{Git Integration:} We will integrate directly with GitHub and GitLab to manage version history seamlessly within the UI.
\end{itemize}